\documentclass[aspectratio=169]{beamer}
\usetheme{Madrid}
\usecolortheme{default}
\usepackage[utf8]{inputenc}
\usepackage[spanish]{babel}
\usepackage{graphicx}
\usepackage{hyperref}

\title{Normas APA - Séptima Edición}
\subtitle{Estructura del Trabajo Académico}
\author{Edison Achalma}
\institute{Universidad Nacional de San Cristóbal de Huamanga}
\date{\today}

\begin{document}

\frame{\titlepage}

\begin{frame}{Tabla de Contenido}
\tableofcontents
\end{frame}

\section{Orden de los Elementos}

\begin{frame}{Elementos del Trabajo Académico}
\begin{enumerate}
    \item Página de presentación (portada)
    \item Resumen
    \item Palabras clave
    \item Contenido o cuerpo del texto
    \item Referencias
    \item Notas al pie (opcional)
    \item Tablas (opcional)
    \item Figuras (opcional)
    \item Anexos (opcional)
\end{enumerate}
\end{frame}

\begin{frame}{Esquema General}
% Imagen: Diagrama de flujo mostrando el orden de los elementos
\begin{center}
\includegraphics[width=0.85\textwidth]{images/esquema_estructura.png}
\end{center}
\end{frame}

\section{Página de Presentación}

\begin{frame}{Portada para Trabajos Estudiantiles}
\textbf{Elementos requeridos (de arriba hacia abajo):}
\begin{enumerate}
    \item Título del trabajo
    \item Nombre del autor(a)
    \item Afiliación (Departamento, Universidad)
    \item Nombre de la asignatura
    \item Nombre del profesor(a)
    \item Fecha de entrega
\end{enumerate}
\end{frame}

\begin{frame}{Ejemplo de Portada}
% Imagen: Ejemplo completo de una portada estudiantil según APA 7
\begin{center}
\includegraphics[width=0.65\textwidth]{images/ejemplo_portada_estudiante.png}
\end{center}
\end{frame}

\begin{frame}{Detalles de la Portada}
\begin{block}{Importante}
\begin{itemize}
    \item Todo el texto va centrado
    \item Números de página desde la portada (esquina superior derecha)
    \item Sin negritas ni cursivas (excepto si el título contiene nombres científicos)
    \item Márgenes de 2.54 cm en todos los lados
\end{itemize}
\end{block}

% Imagen: Detalle de márgenes y numeración de la portada
\begin{center}
\includegraphics[width=0.5\textwidth]{images/detalles_portada.png}
\end{center}
\end{frame}

\section{Resumen}

\begin{frame}{El Resumen (Abstract)}
\begin{columns}
\column{0.5\textwidth}
\textbf{Características:}
\begin{itemize}
    \item Extensión: máximo 250 palabras
    \item Descripción breve y completa
    \item Título "Resumen" centrado y en negrita
    \item Sin sangría en la primera línea
    \item Puede ser estructurado o no
\end{itemize}

\column{0.5\textwidth}
% Imagen: Ejemplo de resumen en formato APA
\includegraphics[width=\textwidth]{images/ejemplo_resumen.png}
\end{columns}
\end{frame}

\begin{frame}{Resumen Estructurado}
% Imagen: Ejemplo de resumen estructurado con etiquetas (Objetivo, Método, Resultados, Conclusiones)
\begin{center}
\includegraphics[width=0.8\textwidth]{images/resumen_estructurado.png}
\end{center}
\end{frame}

\section{Palabras Clave}

\begin{frame}{Palabras Clave (Keywords)}
% Imagen: Ejemplo de cómo formatear las palabras clave
\begin{center}
\includegraphics[width=0.75\textwidth]{images/palabras_clave.png}
\end{center}

\begin{block}{Formato}
\begin{itemize}
    \item Sangría de 1.27 cm
    \item "Palabras clave:" en cursiva
    \item Palabras separadas por comas
    \item En minúscula (excepto nombres propios)
\end{itemize}
\end{block}
\end{frame}

\section{Cuerpo del Texto}

\begin{frame}{Estructura del Contenido}
\textbf{Estructura común (IMRD):}
\begin{itemize}
    \item \textbf{I}ntroducción
    \item \textbf{M}étodo
    \item \textbf{R}esultados
    \item \textbf{D}iscusión
\end{itemize}

\vspace{0.5cm}
% Imagen: Diagrama IMRD con descripción de cada sección
\begin{center}
\includegraphics[width=0.7\textwidth]{images/estructura_imrd.png}
\end{center}
\end{frame}

\begin{frame}{Niveles de Títulos}
% Imagen: Tabla completa mostrando los 5 niveles de títulos con ejemplos
\begin{center}
\includegraphics[width=0.95\textwidth]{images/niveles_titulos.png}
\end{center}
\end{frame}

\begin{frame}{Ejemplos de Niveles de Títulos}
% Imagen: Ejemplo práctico de un texto con diferentes niveles de títulos aplicados
\begin{center}
\includegraphics[width=0.85\textwidth]{images/ejemplo_niveles_titulos.png}
\end{center}
\end{frame}

\section{Referencias}

\begin{frame}{Apartado de Referencias}
\begin{columns}
\column{0.5\textwidth}
\textbf{Características:}
\begin{itemize}
    \item Título "Referencias" centrado y en negrita
    \item Orden alfabético por apellido
    \item Sangría francesa de 1.27 cm
    \item Interlineado doble
    \item Nueva página
\end{itemize}

\column{0.5\textwidth}
% Imagen: Ejemplo de página de referencias formateada
\includegraphics[width=\textwidth]{images/ejemplo_referencias.png}
\end{columns}
\end{frame}

\section{Notas al Pie}

\begin{frame}{Notas al Pie}
% Imagen: Ejemplo de texto con nota al pie correctamente formateada
\begin{center}
\includegraphics[width=0.8\textwidth]{images/ejemplo_notas_pie.png}
\end{center}

\begin{block}{Uso}
Para información complementaria que distraería en el texto principal
\end{block}
\end{frame}

\section{Anexos}

\begin{frame}{Anexos}
\textbf{Contenido de los anexos:}
\begin{itemize}
    \item Materiales utilizados en el estudio
    \item Instrumentos de medición
    \item Cuestionarios
    \item Datos complementarios
\end{itemize}

\vspace{0.5cm}
% Imagen: Ejemplo de portada de un anexo (Anexo A, título descriptivo)
\begin{center}
\includegraphics[width=0.6\textwidth]{images/ejemplo_anexo.png}
\end{center}
\end{frame}

\begin{frame}{Etiquetado de Anexos}
\begin{block}{Formato}
\begin{itemize}
    \item Etiqueta: Anexo A, Anexo B, Anexo C...
    \item Cada anexo en página separada
    \item Título descriptivo bajo la etiqueta
    \item Mencionar explícitamente en el texto: (ver Anexo A)
\end{itemize}
\end{block}

% Imagen: Ejemplo de referencia a anexos en el texto
\begin{center}
\includegraphics[width=0.65\textwidth]{images/referencia_anexos.png}
\end{center}
\end{frame}

\begin{frame}{Diferencia: Referencias vs Bibliografía}
% Imagen: Tabla comparativa entre Referencias y Bibliografía
\begin{center}
\includegraphics[width=0.9\textwidth]{images/referencias_vs_bibliografia.png}
\end{center}

\begin{alertblock}{En APA usamos REFERENCIAS}
Solo se incluyen las fuentes que fueron citadas en el trabajo
\end{alertblock}
\end{frame}

\begin{frame}[standout]
\Huge ¿Preguntas?
\end{frame}

\begin{frame}{Próxima Clase}
\begin{center}
\Large \textbf{Formato General del Documento}
\vspace{1cm}

Veremos:
\begin{itemize}
    \item Márgenes y espaciado
    \item Fuentes tipográficas
    \item Alineación y sangría
    \item Numeración de páginas
\end{itemize}
\end{center}
\end{frame}

\end{document}
